\scanpage{7}

\mspage{444} \handwritten{\unsure{h|4}}

d'avance (qu'il qualifiera de ``vraie'', ou de ``vérité'') n'aura aucune peine à la trouver, et même à la prouver à son entière satisfaction - et sûrement il trouvera bien en chemin tel ou tel autre, si ce n'est toute une foule, tout content de conclure alliance avec lui et de partager convictions et satisfaction. Il est comme le chasseur de papillons, qui part avec dans son filet un beau papillon (empaillé si ça se trouve), et qui le sort tout content ( et à son entière satisfaction) en revenant de sa ``chasse''.

Et il y a celui aussi qui se trouve placé devant un inconnu, comme un enfant nu devant la mer. Quand l'enfant désire la connaître, il entre et la connaît - qu'elle soit tiède ou fraîche, calme ou agitée. Celui qu'attire telle chose inconnue, et qui part pour la connaître, sûrement la connaîtra peu ou prou\unsure{.} Avec ou sans filet, il trouvera le vrai, ou en tous cas. \mbox{d u}\quad vrai. Ses erreur\unsure{s} comme ses trouvailles sont autant d'étapes dans son cheminement, ou pour mieux dire, dans \mbox{s e s}\quad \mbox{a m o u r s} avec ce qu'il désire connaître.

Je sais bien de quoi je parle, car dans ma vie j'ai été abondamment tour à tour, et ce chasseur de papillons, et cet enfant nu. Il n'y a aucune difficult\unsure{é} à distinguer l'un de l'autre. Je doute que les ``critères objectifs'' soient ici d'un grand secours, c'est beaucoup plus simple que çà! Il n'y a qu'à se servir de ses yeux\ldots

Et il n'y a aucune difficulté non plus à distinguer les étapes succéssives\unsure{,} les stades de décantation successifs, dans ce cheminement dont je viens de parle\unsure{r} à partir de cette étape ``morte'' ou nul pressenti \deleted{\unsure{e}}\handwritten{a}ffleurant à la conscience ne fait encore soupçonner ``quelque chose'', au delà d'une certaine surface plate et amorphe que nous présentent des yeux somnolents, et qui à travers des ``éveils'' successifs nous amène vers une appréhension de plus en plus délicate, plus intime, plus complète de ce ``quelque chose''. Il n'est pas de nature essentiellemen\unsure{t} différente, qu'il s'agisse du cheminement dans la découverte des choses mathématiques, ou dans celle de soi et d'autrui. Le sentiment d'une \mbox{p r o g r e s s i o n} dans une \mbox{c o n n a i s s a n c e}, qui \deleted{\unsure{a}}\handwritten{s}'approfondit peu à peu (fût-ce à travers une accumulation d'erreurs, patiemment, inlassablement corrigées) - ce sentiment est aussi irrécusable dans ce dernier cas comme dans l'autre.

Cette \mbox{a s s u r a n c e} - là est l'une des faces d'une disposition intérieure, dont l'autre face est une \mbox{o u v e r t u r e}\quad \mbox{a u}\quad \mbox{d o u t e}: une attitude de curiosité excluant toute crainte, vis à vis de ses propres erreurs, qui permet de les dépister et de les corriger constamment. La condition essentie\unsure{l}le de cette double assise, de cette \mbox{f o i} indispensable pour accueillir le doute comme pour découvrir, est l'absence de toute peur (qu'elle soit apparente
% archive sheet number: 6
